\documentclass[12pt,a4paper]{article}
\usepackage[utf8]{inputenc}
\usepackage[brazil]{babel}
\usepackage[T1]{fontenc}
\usepackage{geometry}
\usepackage{graphicx}
\usepackage{hyperref}
\usepackage{listings}
\usepackage{xcolor}
\usepackage{amsmath}
\usepackage{amssymb}
\usepackage{mathpartir}

\geometry{a4paper, left=3cm, right=2cm, top=3cm, bottom=2cm}

\definecolor{codegreen}{rgb}{0,0.6,0}
\definecolor{codegray}{rgb}{0.5,0.5,0.5}
\definecolor{codepurple}{rgb}{0.58,0,0.82}
\definecolor{backcolour}{rgb}{0.95,0.95,0.92}

\lstdefinestyle{mystyle}{
    backgroundcolor=\color{backcolour},
    commentstyle=\color{codegreen},
    keywordstyle=\color{magenta},
    numberstyle=\tiny\color{codegray},
    stringstyle=\color{codepurple},
    basicstyle=\ttfamily\footnotesize,
    breakatwhitespace=false,
    breaklines=true,
    captionpos=b,
    keepspaces=true,
    numbers=left,
    numbersep=5pt,
    showspaces=false,
    showstringspaces=false,
    showtabs=false,
    tabsize=2
}

\lstdefinestyle{haskell}{
    backgroundcolor=\color{backcolour},
    commentstyle=\color{codegreen},
    keywordstyle=\color{blue},
    numberstyle=\tiny\color{codegray},
    stringstyle=\color{codepurple},
    basicstyle=\ttfamily\footnotesize,
    breakatwhitespace=false,
    breaklines=true,
    captionpos=b,
    keepspaces=true,
    numbers=left,
    numbersep=5pt,
    showspaces=false,
    showstringspaces=false,
    showtabs=false,
    tabsize=2,
    morekeywords={module,where,import,data,deriving,instance,class,type,newtype,let,in,case,of,if,then,else,do,return,qualified,as}
}

\lstdefinestyle{sl}{
    backgroundcolor=\color{backcolour},
    commentstyle=\color{codegreen},
    keywordstyle=\color{blue},
    numberstyle=\tiny\color{codegray},
    stringstyle=\color{codepurple},
    basicstyle=\ttfamily\footnotesize,
    breakatwhitespace=false,
    breaklines=true,
    captionpos=b,
    keepspaces=true,
    numbers=left,
    numbersep=5pt,
    showspaces=false,
    showstringspaces=false,
    showtabs=false,
    tabsize=2,
    morekeywords={func,struct,if,else,while,for,return,let,new,forall,int,float,string,bool,void,true,false,print}
}

\lstset{style=mystyle}

\title{Relatório de Projeto: Compilador SL\\Etapa 2 -- Análise Semântica e Interpretador}
\author{
    Kayo Xavier -- Matrícula: 21.2.4095 \\
    Pedro Vieira -- Matrícula: 20.2.4008 \\[0.5em]
    BCC328 -- Construção de Compiladores I -- DECOM/UFOP
}
\date{\today}

\begin{document}

\maketitle

\begin{abstract}
Este relatório descreve a segunda etapa do desenvolvimento do compilador para a linguagem SL (\textit{Simple Language}), implementado em Haskell. Nesta etapa, foram projetados e implementados o analisador semântico e o interpretador da linguagem. O analisador semântico realiza verificações de tipos, resolução de nomes e validação de escopos, enquanto o interpretador executa programas SL por meio de caminhamento na árvore de sintaxe abstrata (AST). O documento apresenta o sistema de tipos formalizado, a semântica operacional de \textit{big-step}, as decisões de projeto e os resultados obtidos nos testes com os programas de exemplo da especificação.
\end{abstract}

\tableofcontents

\newpage

%% ================================================================
\section{Introdução}
%% ================================================================

A linguagem SL (\textit{Simple Language}) é uma linguagem de programação imperativa com tipagem estática, projetada para fins didáticos no contexto da disciplina BCC328 -- Construção de Compiladores I. A linguagem possui suporte a tipos primitivos (\texttt{int}, \texttt{float}, \texttt{string}, \texttt{bool}), estruturas de controle (\texttt{if-else}, \texttt{while}, \texttt{for}), arranjos unidimensionais, registros (\textit{structs}), funções com parâmetros e retorno, polimorfismo paramétrico (\texttt{forall}) e inferência de tipos.

Na primeira etapa do trabalho, foram implementados o analisador léxico e o analisador sintático utilizando a biblioteca Megaparsec, resultando na construção da árvore de sintaxe abstrata (AST) a partir do código-fonte SL.

Nesta segunda etapa, o objetivo foi implementar duas componentes fundamentais:

\begin{enumerate}
    \item \textbf{Analisador Semântico}: responsável por verificar a correção semântica dos programas, incluindo compatibilidade de tipos, unicidade de declarações, verificação de chamadas de função e acesso válido a campos de registros e elementos de arranjos.
    \item \textbf{Interpretador}: responsável por executar programas SL diretamente a partir da AST, implementando a semântica operacional da linguagem.
\end{enumerate}

O compilador foi desenvolvido inteiramente em Haskell, utilizando o sistema de build Cabal, e está preparado para execução em ambiente Docker conforme especificado pelo professor.

%% ================================================================
\section{Metodologia}
%% ================================================================

\subsection{Estrutura sintática de SL}

A gramática de SL, definida na Etapa 1, segue uma estrutura hierárquica composta por declarações de topo (\textit{top-level}), que podem ser definições de funções ou de registros. Cada função possui um nome, uma lista de parâmetros (opcionalmente tipados), um tipo de retorno opcional e um corpo composto por instruções (\textit{statements}). A gramática é processada por um analisador sintático descendente recursivo implementado com Megaparsec e a biblioteca \texttt{parser-combinators} para tratamento de precedência de operadores.

A AST resultante é definida no módulo \texttt{SL.AST} com os seguintes tipos principais:

\begin{itemize}
    \item \texttt{Program}: lista de declarações de topo (\texttt{TopLevel})
    \item \texttt{TopLevel}: \texttt{TLFunc} (funções) ou \texttt{TLStruct} (registros)
    \item \texttt{Stmt}: instruções como \texttt{SLet}, \texttt{SAssign}, \texttt{SIf}, \texttt{SWhile}, \texttt{SFor}, \texttt{SReturn}, \texttt{SBlock}
    \item \texttt{Expr}: expressões como literais, variáveis, operações binárias/unárias, chamadas de função, acesso a arranjos/campos, construtores de arranjo e registro
    \item \texttt{Type}: tipos como \texttt{TyInt}, \texttt{TyFloat}, \texttt{TyString}, \texttt{TyBool}, \texttt{TyVoid}, \texttt{TyVar} (variável de tipo), \texttt{TyArray}, \texttt{TyStruct}, \texttt{TyFunc}
\end{itemize}

%% ----------------------------------------------------------------
\subsection{Sistema de tipos para SL}
%% ----------------------------------------------------------------

O sistema de tipos de SL é estático e nominalmente tipado. Definimos formalmente os tipos da linguagem e as regras de tipagem utilizadas pelo analisador semântico.

\subsubsection{Tipos da linguagem}

Os tipos de SL são definidos pela seguinte gramática:

\[
\tau ::= \texttt{int} \mid \texttt{float} \mid \texttt{string} \mid \texttt{bool} \mid \texttt{void} \mid \alpha \mid \tau[\,] \mid S \mid (\tau_1, \ldots, \tau_n) \to \tau
\]

\noindent onde $\alpha$ denota uma variável de tipo (utilizada em funções genéricas), $\tau[\,]$ denota um arranjo de elementos do tipo $\tau$, $S$ denota um tipo de registro (struct) definido pelo usuário, e $(\tau_1, \ldots, \tau_n) \to \tau$ denota um tipo funcional.

\subsubsection{Compatibilidade de tipos}

Dois tipos $\tau_1$ e $\tau_2$ são considerados compatíveis ($\tau_1 \sim \tau_2$) se satisfizerem uma das seguintes condições:

\begin{enumerate}
    \item $\tau_1 = \tau_2$ (igualdade estrutural)
    \item $\tau_1$ ou $\tau_2$ é uma variável de tipo $\alpha$ (unificação com qualquer tipo)
    \item $\tau_1, \tau_2 \in \{\texttt{int}, \texttt{float}\}$ (coerção numérica implícita)
    \item $\tau_1 = \tau'_1[\,]$ e $\tau_2 = \tau'_2[\,]$ com $\tau'_1 \sim \tau'_2$ (compatibilidade de arranjos)
    \item $\tau_1 = (\sigma_1, \ldots, \sigma_n) \to \rho_1$ e $\tau_2 = (\sigma'_1, \ldots, \sigma'_n) \to \rho_2$ com $\sigma_i \sim \sigma'_i$ para todo $i$ e $\rho_1 \sim \rho_2$ (compatibilidade de tipos funcionais)
\end{enumerate}

\subsubsection{Regras de tipagem}

As principais regras de tipagem são apresentadas a seguir, usando a notação $\Gamma \vdash e : \tau$ para indicar que no ambiente $\Gamma$, a expressão $e$ possui tipo $\tau$.

\paragraph{Literais.}

\begin{mathpar}
\inferrule{ }{\Gamma \vdash n : \texttt{int}} \and
\inferrule{ }{\Gamma \vdash f : \texttt{float}} \and
\inferrule{ }{\Gamma \vdash s : \texttt{string}} \and
\inferrule{ }{\Gamma \vdash b : \texttt{bool}}
\end{mathpar}

\paragraph{Variáveis.}

\begin{mathpar}
\inferrule{x : \tau \in \Gamma}{\Gamma \vdash x : \tau}
\end{mathpar}

\paragraph{Operações aritméticas.}

\begin{mathpar}
\inferrule{\Gamma \vdash e_1 : \tau_1 \\ \Gamma \vdash e_2 : \tau_2 \\ \tau_1, \tau_2 \in \{\texttt{int}, \texttt{float}\} \\ \oplus \in \{+, -, *, /, \%\}}
          {\Gamma \vdash e_1 \oplus e_2 : \text{wider}(\tau_1, \tau_2)}
\end{mathpar}

\noindent onde $\text{wider}(\tau_1, \tau_2) = \texttt{float}$ se $\tau_1 = \texttt{float}$ ou $\tau_2 = \texttt{float}$; caso contrário, $\text{wider}(\tau_1, \tau_2) = \texttt{int}$.

\paragraph{Operações de comparação.}

\begin{mathpar}
\inferrule{\Gamma \vdash e_1 : \tau_1 \\ \Gamma \vdash e_2 : \tau_2 \\ \tau_1, \tau_2 \in \{\texttt{int}, \texttt{float}\} \\ \oplus \in \{<, \leq, >, \geq\}}
          {\Gamma \vdash e_1 \oplus e_2 : \texttt{bool}}
\end{mathpar}

\paragraph{Operações de igualdade.}

\begin{mathpar}
\inferrule{\Gamma \vdash e_1 : \tau_1 \\ \Gamma \vdash e_2 : \tau_2 \\ \tau_1 \sim \tau_2 \\ \oplus \in \{==, !=\}}
          {\Gamma \vdash e_1 \oplus e_2 : \texttt{bool}}
\end{mathpar}

\paragraph{Operações lógicas.}

\begin{mathpar}
\inferrule{\Gamma \vdash e_1 : \texttt{bool} \\ \Gamma \vdash e_2 : \texttt{bool} \\ \oplus \in \{\&\&, ||\}}
          {\Gamma \vdash e_1 \oplus e_2 : \texttt{bool}}
\end{mathpar}

\paragraph{Operações unárias.}

\begin{mathpar}
\inferrule{\Gamma \vdash e : \tau \\ \tau \in \{\texttt{int}, \texttt{float}\}}
          {\Gamma \vdash -e : \tau}
\and
\inferrule{\Gamma \vdash e : \texttt{bool}}
          {\Gamma \vdash !e : \texttt{bool}}
\end{mathpar}

\paragraph{Chamada de função.}

\begin{mathpar}
\inferrule{f : (\tau_1, \ldots, \tau_n) \to \tau_r \in \Gamma \\ \Gamma \vdash a_i : \sigma_i \\ \tau_i \sim \sigma_i \text{ para } 1 \leq i \leq n}
          {\Gamma \vdash f(a_1, \ldots, a_n) : \tau_r}
\end{mathpar}

\paragraph{Acesso a arranjo.}

\begin{mathpar}
\inferrule{\Gamma \vdash e_1 : \tau[\,] \\ \Gamma \vdash e_2 : \texttt{int}}
          {\Gamma \vdash e_1[e_2] : \tau}
\end{mathpar}

\paragraph{Acesso a campo de registro.}

\begin{mathpar}
\inferrule{\Gamma \vdash e : S \\ \text{field}(S, f) = \tau}
          {\Gamma \vdash e.f : \tau}
\end{mathpar}

\paragraph{Propriedade \texttt{.size} de arranjos.}

\begin{mathpar}
\inferrule{\Gamma \vdash e : \tau[\,]}
          {\Gamma \vdash e.\texttt{size} : \texttt{int}}
\end{mathpar}

%% ----------------------------------------------------------------
\subsection{Inferência de tipos para SL}
%% ----------------------------------------------------------------

A inferência de tipos em SL é utilizada em dois contextos principais:

\begin{enumerate}
    \item \textbf{Variáveis sem anotação de tipo}: quando uma declaração \texttt{let x = e;} não possui anotação de tipo, o tipo de \texttt{x} é inferido a partir do tipo da expressão \texttt{e}.
    \item \textbf{Parâmetros de função sem anotação}: quando um parâmetro de função não possui anotação de tipo (como em \texttt{func id(x) \{ return x; \}}), uma variável de tipo fresca $\_t_i$ é gerada e associada ao parâmetro.
\end{enumerate}

A geração de variáveis de tipo frescas é implementada por um contador monotonicamente crescente no ambiente de checagem. Cada nova variável recebe o nome $\_t_0, \_t_1, \ldots$, garantindo unicidade.

Para funções genéricas declaradas com \texttt{forall}, as variáveis de tipo explícitas (e.g., \texttt{a}, \texttt{b}) são registradas no ambiente como \texttt{TyVar}, permitindo que sejam unificadas com qualquer tipo concreto durante a verificação.

%% ----------------------------------------------------------------
\subsection{Semântica operacional para SL}
%% ----------------------------------------------------------------

A semântica operacional de SL é definida no estilo \textit{big-step} (semântica natural), onde a avaliação de uma expressão ou execução de uma instrução produz diretamente o resultado final. A notação utilizada é:

\begin{itemize}
    \item $\langle e, \sigma \rangle \Downarrow v$ : a expressão $e$ avaliada no estado $\sigma$ produz o valor $v$
    \item $\langle s, \sigma \rangle \Downarrow \sigma'$ : a instrução $s$ executada no estado $\sigma$ produz um novo estado $\sigma'$
\end{itemize}

\noindent onde $\sigma$ é um ambiente de execução que mapeia nomes a valores.

\subsubsection{Valores de tempo de execução}

Os valores de tempo de execução são definidos por:

\[
v ::= n \mid f \mid s \mid b \mid \texttt{void} \mid \langle\texttt{arr}, \text{ref}\rangle \mid \langle\texttt{struct}, S, \text{campos}\rangle \mid \langle\texttt{func}, \text{def}\rangle
\]

\noindent onde $n$ é um inteiro, $f$ um ponto flutuante, $s$ uma string, $b$ um booleano, $\langle\texttt{arr}, \text{ref}\rangle$ uma referência mutável a um arranjo, e $\langle\texttt{struct}, S, \text{campos}\rangle$ uma referência mutável a um registro.

\subsubsection{Regras para expressões}

\paragraph{Literais.}

\begin{mathpar}
\inferrule{ }{\langle n, \sigma \rangle \Downarrow n}
\and
\inferrule{ }{\langle f, \sigma \rangle \Downarrow f}
\and
\inferrule{ }{\langle s, \sigma \rangle \Downarrow s}
\and
\inferrule{ }{\langle b, \sigma \rangle \Downarrow b}
\end{mathpar}

\paragraph{Variável.}

\begin{mathpar}
\inferrule{\sigma(x) = v}{\langle x, \sigma \rangle \Downarrow v}
\end{mathpar}

\paragraph{Operação binária.}

\begin{mathpar}
\inferrule{\langle e_1, \sigma \rangle \Downarrow v_1 \\ \langle e_2, \sigma \rangle \Downarrow v_2 \\ v_1 \oplus v_2 = v}
          {\langle e_1 \oplus e_2, \sigma \rangle \Downarrow v}
\end{mathpar}

\paragraph{Chamada de função.}

\begin{mathpar}
\inferrule{f \mapsto (\overline{p_i}, \text{corpo}) \in \sigma \\ \langle a_i, \sigma \rangle \Downarrow v_i \\ \sigma' = \sigma[p_i \mapsto v_i] \\ \langle \text{corpo}, \sigma' \rangle \Downarrow_{\text{ret}} v}
          {\langle f(\overline{a_i}), \sigma \rangle \Downarrow v}
\end{mathpar}

\subsubsection{Regras para instruções}

\paragraph{Declaração de variável (\texttt{let}).}

\begin{mathpar}
\inferrule{\langle e, \sigma \rangle \Downarrow v}
          {\langle \texttt{let}\ x\ \texttt{=}\ e, \sigma \rangle \Downarrow \sigma[x \mapsto v]}
\and
\inferrule{\tau = \text{tipo anotado} \\ v_0 = \text{default}(\tau)}
          {\langle \texttt{let}\ x : \tau, \sigma \rangle \Downarrow \sigma[x \mapsto v_0]}
\end{mathpar}

\paragraph{Atribuição.}

\begin{mathpar}
\inferrule{\langle e, \sigma \rangle \Downarrow v}
          {\langle x\ \texttt{=}\ e, \sigma \rangle \Downarrow \sigma[x \mapsto v]}
\end{mathpar}

\paragraph{Condicional (\texttt{if-else}).}

\begin{mathpar}
\inferrule{\langle e, \sigma \rangle \Downarrow \texttt{true} \\ \langle s_1, \sigma \rangle \Downarrow \sigma'}
          {\langle \texttt{if}\ (e)\ s_1\ \texttt{else}\ s_2, \sigma \rangle \Downarrow \sigma'}
\and
\inferrule{\langle e, \sigma \rangle \Downarrow \texttt{false} \\ \langle s_2, \sigma \rangle \Downarrow \sigma'}
          {\langle \texttt{if}\ (e)\ s_1\ \texttt{else}\ s_2, \sigma \rangle \Downarrow \sigma'}
\end{mathpar}

\paragraph{Laço \texttt{while}.}

\begin{mathpar}
\inferrule{\langle e, \sigma \rangle \Downarrow \texttt{true} \\ \langle s, \sigma \rangle \Downarrow \sigma' \\ \langle \texttt{while}\ (e)\ s, \sigma' \rangle \Downarrow \sigma''}
          {\langle \texttt{while}\ (e)\ s, \sigma \rangle \Downarrow \sigma''}
\and
\inferrule{\langle e, \sigma \rangle \Downarrow \texttt{false}}
          {\langle \texttt{while}\ (e)\ s, \sigma \rangle \Downarrow \sigma}
\end{mathpar}

\paragraph{Retorno.}

\begin{mathpar}
\inferrule{\langle e, \sigma \rangle \Downarrow v}
          {\langle \texttt{return}\ e, \sigma \rangle \Downarrow_{\text{ret}} v}
\end{mathpar}

%% ================================================================
\section{Arquitetura do Compilador}
%% ================================================================

O compilador SL é organizado em uma arquitetura modular, com cada fase implementada como um módulo Haskell independente. A estrutura segue o padrão clássico de \textit{pipeline} de compilação:

\begin{center}
\texttt{Código-fonte} $\xrightarrow{\text{Lexer}}$ \texttt{Tokens} $\xrightarrow{\text{Parser}}$ \texttt{AST} $\xrightarrow{\text{TypeChecker}}$ \texttt{AST validada} $\xrightarrow{\text{Interpreter}}$ \texttt{Execução}
\end{center}

%% ----------------------------------------------------------------
\subsection{Análise léxica}
%% ----------------------------------------------------------------

O analisador léxico é implementado no módulo \texttt{SL.Lexer} utilizando a biblioteca Megaparsec. O lexer converte o código-fonte em uma sequência de \textit{tokens}, cada um acompanhado de informação de posição (linha e coluna).

Os tokens reconhecidos incluem:
\begin{itemize}
    \item Literais: inteiros, ponto flutuante, strings e booleanos
    \item Identificadores e palavras reservadas: \texttt{func}, \texttt{struct}, \texttt{if}, \texttt{else}, \texttt{while}, \texttt{for}, \texttt{return}, \texttt{let}, \texttt{new}, \texttt{forall}, e tipos (\texttt{int}, \texttt{float}, \texttt{string}, \texttt{bool}, \texttt{void})
    \item Operadores: \texttt{+}, \texttt{-}, \texttt{*}, \texttt{/}, \texttt{\%}, \texttt{==}, \texttt{!=}, \texttt{<}, \texttt{<=}, \texttt{>}, \texttt{>=}, \texttt{\&\&}, \texttt{||}, \texttt{!}
    \item Pontuação: parênteses, chaves, colchetes, ponto-e-vírgula, vírgula, dois-pontos, ponto, seta (\texttt{->})
\end{itemize}

O lexer também trata comentários de linha (\texttt{//}) e de bloco (\texttt{/* */}).

%% ----------------------------------------------------------------
\subsection{Análise sintática}
%% ----------------------------------------------------------------

O analisador sintático é implementado no módulo \texttt{SL.Parser}, também utilizando Megaparsec. O parser é um analisador descendente recursivo que constrói a AST diretamente.

A precedência de operadores é tratada pela função \texttt{makeExprParser} da biblioteca \texttt{parser-combinators}, com a seguinte tabela de precedência (da mais alta para a mais baixa):

\begin{enumerate}
    \item Operadores unários: \texttt{-} (negação), \texttt{!} (negação lógica)
    \item Multiplicativos: \texttt{*}, \texttt{/}, \texttt{\%}
    \item Aditivos: \texttt{+}, \texttt{-}
    \item Comparação: \texttt{<=}, \texttt{>=}, \texttt{<}, \texttt{>}
    \item Igualdade: \texttt{==}, \texttt{!=}
    \item E lógico: \texttt{\&\&}
    \item Ou lógico: \texttt{||}
\end{enumerate}

Expressões pós-fixas (acesso a arranjo \texttt{[i]}, acesso a campo \texttt{.field} e chamada de função \texttt{(args)}) são tratadas como operações de \textit{postfix} aplicadas ao termo base.

%% ----------------------------------------------------------------
\subsection{Árvore de sintaxe abstrata}
%% ----------------------------------------------------------------

A AST é definida no módulo \texttt{SL.AST} utilizando tipos algébricos de Haskell. As principais decisões de projeto incluem:

\begin{itemize}
    \item \textbf{Tipos opcionais em parâmetros}: \texttt{FuncParam} possui \texttt{paramType :: Maybe Type}, permitindo que parâmetros sem anotação tenham seu tipo inferido.
    \item \textbf{Tipos de retorno opcionais}: \texttt{funcRetType :: Maybe Type} em \texttt{TLFunc} permite inferência do tipo de retorno.
    \item \textbf{Variáveis de tipo genérico}: \texttt{funcTypeVars :: [TypeVar]} armazena as variáveis de tipo declaradas com \texttt{forall}.
    \item \textbf{Tipo funcional}: \texttt{TyFunc [Type] Type} representa tipos de funções como \texttt{(a) -> b}, suportando funções de alta ordem.
    \item \textbf{Tamanho opcional em arranjos}: \texttt{TyArray Type (Maybe Expr)} permite tanto \texttt{int[]} quanto \texttt{int[5]}.
\end{itemize}

%% ----------------------------------------------------------------
\subsection{Análise semântica}
%% ----------------------------------------------------------------

O analisador semântico é implementado nos módulos \texttt{SL.TypeChecker} e \texttt{SL.Env}. A verificação é realizada em duas passadas sobre a AST:

\paragraph{Passada 1 -- Registro de declarações:} Todas as funções e structs são registradas no ambiente global antes de verificar os corpos. Isso permite referências mútuas entre funções e o uso de structs antes de suas definições.

\paragraph{Passada 2 -- Verificação de corpos:} Cada corpo de função é verificado quanto a:

\begin{itemize}
    \item \textbf{Resolução de nomes}: verifica se variáveis e funções referenciadas estão declaradas no escopo visível.
    \item \textbf{Compatibilidade de tipos}: valida operações aritméticas, lógicas, de comparação e de igualdade conforme as regras de tipagem.
    \item \textbf{Declaração única no escopo}: impede a redeclaração de variáveis no mesmo escopo utilizando \texttt{lookupVarCurrentScope}.
    \item \textbf{Verificação de chamadas de função}: confere o número e tipos dos argumentos contra a assinatura registrada.
    \item \textbf{Acesso a campos}: valida que acessos \texttt{e.field} referem-se a campos existentes do struct correspondente.
    \item \textbf{Indexação de arranjos}: verifica que o operando é um arranjo e que o índice é do tipo \texttt{int}.
    \item \textbf{Tipos de retorno}: valida que instruções \texttt{return} são compatíveis com o tipo de retorno declarado.
\end{itemize}

O ambiente de verificação (\texttt{SL.Env}) utiliza uma pilha de escopos (lista de \texttt{Map Text Type}) para variáveis, com mapas globais para funções e structs. Erros semânticos são acumulados (não interrompem a verificação), permitindo reportar todos os erros de uma só vez.

\paragraph{Suporte a funções de alta ordem.} Para chamadas como \texttt{f(v[i])} onde \texttt{f} é um parâmetro de função com tipo \texttt{(a) -> b}, o verificador consulta o ambiente de variáveis caso a função não seja encontrada no registro global, tratando variáveis do tipo \texttt{TyFunc} como funções invocáveis.

%% ----------------------------------------------------------------
\subsection{Interpretador}
%% ----------------------------------------------------------------

O interpretador é implementado no módulo \texttt{SL.Interpreter} como um \textit{tree-walking interpreter} que avalia a AST diretamente na mônada \texttt{IO} de Haskell. As principais decisões de projeto são:

\paragraph{Valores de tempo de execução.} Os valores são representados pelo tipo algébrico \texttt{Value}:

\begin{lstlisting}[style=haskell]
data Value
    = VInt Integer
    | VFloat Double
    | VString Text
    | VBool Bool
    | VArray (IORef [IORef Value])
    | VStruct Text (IORef (Map Text (IORef Value)))
    | VFunc FuncDef
    | VVoid
\end{lstlisting}

\paragraph{Mutabilidade via \texttt{IORef}.} Arranjos e registros são representados usando \texttt{IORef} (referências mutáveis), permitindo que atribuições como \texttt{arr[i] = v} e \texttt{person.name = "Bob"} modifiquem os valores in-place, conforme esperado em uma linguagem imperativa.

\paragraph{Ambiente de execução.} O ambiente de tempo de execução utiliza uma pilha de escopos similar ao analisador semântico, porém mapeando nomes a \texttt{IORef Value} em vez de tipos:

\begin{lstlisting}[style=haskell]
data Env = Env
    { envScopes     :: [Map Text (IORef Value)]
    , envGlobals    :: IORef (Map Text FuncDef)
    , envStructDefs :: Map Text StructDefRT
    }
\end{lstlisting}

\paragraph{Implementação de \texttt{return}.} Instruções \texttt{return} são implementadas utilizando o mecanismo de exceções de Haskell (\texttt{throwIO}/\texttt{catch}). Quando um \texttt{return e} é executado, uma exceção \texttt{ReturnException v} é lançada e capturada pelo \texttt{callFunction} que invocou a função, extraindo o valor de retorno.

\paragraph{Inicialização padrão de variáveis.} Declarações como \texttt{let people : Person[3];} (sem inicializador) criam valores padrão baseados no tipo anotado: \texttt{0} para \texttt{int}, \texttt{0.0} para \texttt{float}, \texttt{""} para \texttt{string}, \texttt{false} para \texttt{bool}, e arranjos pré-alocados com o tamanho especificado para tipos de arranjo.

\paragraph{Execução em três passadas.}
\begin{enumerate}
    \item \textbf{Registro de structs}: todas as definições de struct são coletadas no ambiente.
    \item \textbf{Registro de funções}: todas as definições de função são registradas como \texttt{FuncDef} no mapa global.
    \item \textbf{Chamada de \texttt{main}}: a função \texttt{main} é localizada e invocada com argumentos vazios.
\end{enumerate}

\paragraph{Função \texttt{print} embutida.} A função \texttt{print} é tratada como uma função \textit{built-in} que aceita qualquer número de argumentos e imprime cada valor em uma linha separada, convertendo para sua representação textual.

%% ================================================================
\section{Resultados e Discussão}
%% ================================================================

\subsection{Instruções de Uso}

O compilador pode ser compilado e executado utilizando Cabal dentro do ambiente Docker fornecido:

\begin{lstlisting}[language=bash,style=mystyle]
# Compilar o projeto
cabal build

# Executar um programa SL
cabal run sl-compiler -- --interpret programa.sl

# Modos disponiveis
cabal run sl-compiler -- --lexer programa.sl      # tokens
cabal run sl-compiler -- --parser programa.sl     # AST
cabal run sl-compiler -- --pretty programa.sl     # codigo reconstruido
cabal run sl-compiler -- --typecheck programa.sl  # analise semantica
cabal run sl-compiler -- --interpret programa.sl  # execucao
cabal run sl-compiler -- --run programa.sl        # alias para interpret
\end{lstlisting}

Para iniciar o ambiente Docker:

\begin{lstlisting}[language=bash,style=mystyle]
docker run -it --rm -v ./workspace:/workspace \
    -w /workspace trabalho-02-awsome-team-name-sl bash
\end{lstlisting}

%% ----------------------------------------------------------------
\subsection{Testes Realizados}
%% ----------------------------------------------------------------

O projeto conta com 144 testes automatizados utilizando a biblioteca HSpec, organizados nas seguintes categorias:

\begin{table}[h!]
\centering
\begin{tabular}{|l|c|}
\hline
\textbf{Categoria} & \textbf{Quantidade} \\
\hline
Lexer & 8 \\
Parser & 22 \\
Pretty Printer & 6 \\
Testes Negativos (sintático) & 5 \\
Testes de Integração (sintático) & 8 \\
TypeChecker -- Programas Válidos & 24 \\
TypeChecker -- Programas Inválidos & 21 \\
TypeChecker -- Inferência & 8 \\
TypeChecker -- Integração & 8 \\
Interpreter -- Básico & 9 \\
Interpreter -- Fluxo de Controle & 5 \\
Interpreter -- Funções & 4 \\
Interpreter -- Estruturas de Dados & 5 \\
Interpreter -- Integração & 4 \\
Interpreter -- Erros de Execução & 5 \\
\hline
\textbf{Total} & \textbf{144} \\
\hline
\end{tabular}
\caption{Distribuição dos testes automatizados por categoria.}
\end{table}

\paragraph{Exemplos da especificação.} Todos os programas de exemplo fornecidos na especificação foram testados com sucesso:

\begin{table}[h!]
\centering
\begin{tabular}{|l|l|l|}
\hline
\textbf{Exemplo} & \textbf{Modo} & \textbf{Saída} \\
\hline
1 -- Fatorial & \texttt{--interpret} & \texttt{120} \\
2 -- Registros e Arranjos & \texttt{--interpret} & Alice/25/1.65, Bob/30/1.8, Charlie/35/1.75 \\
3 -- Operações com Arranjos & \texttt{--interpret} & \texttt{5, 4, 3, 2, 1} \\
4 -- Operações Matemáticas & \texttt{--interpret} & \texttt{23.02..., true, Condicao normal} \\
5 -- Função Identidade & \texttt{--typecheck} & Passa (sem \texttt{main}) \\
6 -- Função \texttt{map} Genérica & \texttt{--typecheck} & Passa (sem \texttt{main}) \\
\hline
\end{tabular}
\caption{Resultados da execução dos exemplos da especificação.}
\end{table}

Para executar todos os testes:

\begin{lstlisting}[language=bash,style=mystyle]
cabal test
\end{lstlisting}

%% ----------------------------------------------------------------
\subsection{Limitações}
%% ----------------------------------------------------------------

As principais limitações conhecidas da implementação atual são:

\begin{enumerate}
    \item \textbf{Inferência de tipos limitada}: a inferência de tipos é local (baseada em expressões de inicialização), sem unificação global ao estilo Hindley-Milner. Variáveis de tipo genérico são tratadas como compatíveis com qualquer tipo, mas não há propagação completa de restrições.

    \item \textbf{Ausência de geração de código}: conforme planejado, a geração de código WebAssembly será implementada na Etapa 3.

    \item \textbf{Sem otimizações}: o interpretador avalia expressões por caminhamento direto na AST, sem otimizações como \textit{constant folding} ou eliminação de código morto.

    \item \textbf{Mensagens de erro sem posição precisa}: embora o analisador semântico acumule erros, as posições de linha e coluna não são propagadas desde a AST em todos os casos.

    \item \textbf{Exemplos 5 e 6 apenas verificáveis}: os exemplos 5 (identidade) e 6 (map genérico) não possuem função \texttt{main}, logo não são executáveis pelo interpretador, apenas verificáveis pelo \textit{type checker}.
\end{enumerate}

%% ================================================================
\section{Conclusão}
%% ================================================================

Nesta segunda etapa do trabalho, implementamos com sucesso o analisador semântico e o interpretador para a linguagem SL. O analisador semântico realiza todas as verificações solicitadas na especificação: compatibilidade de tipos em operações, declaração única de identificadores no mesmo escopo, verificação de chamadas de função com número e tipos corretos de argumentos, e acesso válido a campos de registros e elementos de arranjos.

O interpretador implementa a semântica operacional de SL por meio de caminhamento na AST, suportando todas as construções da linguagem: funções (incluindo recursão), registros mutáveis, arranjos mutáveis, estruturas de controle e a função embutida \texttt{print}. A utilização de \texttt{IORef} para representar valores mutáveis mostrou-se uma decisão eficaz, permitindo a semântica de referência necessária para atribuições em arranjos e campos de registros.

Todos os exemplos fornecidos na especificação foram executados com os resultados esperados, e a suíte de 144 testes automatizados valida o comportamento correto das diversas componentes. O projeto está preparado para a terceira etapa, que envolverá a implementação da inferência de tipos completa e a geração de código WebAssembly.

%% ================================================================
\section{Referências}
%% ================================================================

\begin{thebibliography}{9}
\bibitem{hutton2016}
Hutton, G. (2016). \emph{Programming in Haskell}. Cambridge University Press.

\bibitem{appel1998}
Appel, A. W. (1998). \emph{Modern Compiler Implementation in ML}. Cambridge University Press.

\bibitem{webassembly2023}
WebAssembly Community Group. (2023). \emph{WebAssembly Specification}.

\bibitem{megaparsec}
Kireev, M. (2023). \emph{Megaparsec: Industrial-strength monadic parser combinator library}. Disponível em: \url{https://hackage.haskell.org/package/megaparsec}.

\bibitem{pierce2002}
Pierce, B. C. (2002). \emph{Types and Programming Languages}. MIT Press.
\end{thebibliography}

\end{document}
